%% =============================================================
%%  chapters/06_preprocessing_sync.tex
%% =============================================================
\chapter{Data Preprocessing and Synchronisation}
\label{ch:preprocessing-sync}

\section{Pressure Heatmap Preprocessing}
\label{sec:pressure-preprocessing}

\subsection{Noise Filtering and Normalisation}
\label{sec:pressure-noise-filtering}
Raw pressure heatmaps are first passed through a spatial median filter to
suppress isolated sensor cell dropouts and transient noise spikes, after
which pixel values are normalised to a fixed range using the per-sensor
calibration parameters obtained during the setup described in
Section~\ref{sec:sensor-hardware-seat-setup}. This normalisation step
ensures that synthetic and real heatmaps share a consistent value range
prior to being combined in the mixed training dataset.
%% TODO: Report the exact filter kernel size and normalisation bounds used.

\subsection{Spatial Alignment and Resizing}
\label{sec:pressure-spatial-alignment}
Because the physical sensor grid resolution differs from the target input
resolution expected by the pressure branch described in
Section~\ref{sec:pressure-branch}, heatmaps are resampled onto a fixed
target grid size using bilinear interpolation, and a fixed reference
coordinate frame is applied so that seat regions (seat base, backrest)
occupy consistent pixel regions across all recordings.
%% TODO: State the exact target heatmap resolution used as network input.

\section{Camera Frame Preprocessing}
\label{sec:camera-preprocessing}

\subsection{Image Normalisation and Augmentation}
\label{sec:camera-normalisation-augmentation}
Camera frames are resized to a fixed input resolution and normalised
using per-channel mean and standard deviation statistics computed over
the training split. During training, standard augmentations such as
random brightness and contrast jitter, small rotations, and horizontal
flipping (where anatomically valid for the seating position) are applied
to improve robustness to the lighting and pose variation present in the
target deployment scenarios described in Section~\ref{sec:use-case-definition}.
%% TODO: Confirm final augmentation parameter ranges after initial
%% training experiments.

\subsection{Background Removal and ROI Extraction}
\label{sec:camera-roi-extraction}
A fixed region of interest (ROI) corresponding to each monitored seating
position is cropped from the full camera frame using the known camera
mounting geometry, and static background regions (e.g. door panels,
window area) outside the seat envelope are masked out to reduce
irrelevant visual variation presented to the camera branch.
%% TODO: Document the ROI definition procedure (manual annotation vs.
%% calibration-based projection) once finalised.

\section{Temporal Synchronisation}
\label{sec:temporal-synchronisation}

\subsection{Timestamp Alignment Strategy}
\label{sec:timestamp-alignment}
Camera frames and pressure heatmap samples are each recorded with
hardware or system-clock timestamps, and are aligned by matching each
pressure sample to the camera frame with the closest timestamp within a
fixed maximum tolerance window, discarding samples that fall outside this
window to avoid pairing temporally inconsistent data.
%% TODO: Specify the exact timestamp source (hardware trigger vs. system
%% clock) and the chosen tolerance window in milliseconds.

\subsection{Handling Frame Rate Mismatch}
\label{sec:frame-rate-mismatch}
Because the camera and pressure sensor operate at different native
sampling rates, the higher-rate stream is downsampled to match the
lower-rate stream at inference time, while during dataset construction
all valid synchronised pairs are retained to maximise the number of
training samples available from each recording session.
%% TODO: Report native sampling rates for both modalities and the chosen
%% downsampling strategy.

\subsection{Validation of Synchronisation Quality}
\label{sec:sync-quality-validation}
Synchronisation quality is validated by manually inspecting a sample of
paired camera-pressure recordings for temporal consistency (e.g.
confirming that a visible occupant movement in the camera frame
corresponds to a matching pressure distribution change), and by measuring
the distribution of timestamp offsets across all paired samples in the
dataset.
%% TODO: Report the measured timestamp offset distribution once the
%% synchronisation pipeline is implemented.

\section{Final Data Pipeline Overview}
\label{sec:final-pipeline-overview}
Figure~\ref{fig:preprocessing-pipeline} illustrates the complete
preprocessing pipeline, from raw camera and pressure sensor streams
through per-modality preprocessing and temporal synchronisation to the
production of paired training samples consumed by the fusion architecture
described in Chapter~\ref{ch:fusion-architecture}.

\begin{figure}[H]
\centering
%% TODO: Replace with the actual preprocessing pipeline diagram once
%% created.
\includegraphics[width=0.85\textwidth]{preprocessing_pipeline}
\caption{Overview of the data preprocessing and temporal synchronisation
pipeline, from raw sensor streams to synchronised training samples.}
\label{fig:preprocessing-pipeline}
\end{figure}
